\subsection{1. Mô tả chi tiết bài toán}
\subsubsection{1.1. Bài toán phát hiện Sybil trên Web3}
\indent\indent Nghiên cứu tập trung giải quyết bài toán phát hiện tài khoản giả mạo (Sybil) trên mạng xã hội phi tập trung (Lens Protocol) bằng cách phân loại các thực thể người dùng vào các nhóm rủi ro đã được định nghĩa. Khác với các phương pháp phát hiện Sybil truyền thống trên Web2 (như Facebook hay Twitter) thường chỉ dựa vào hành vi đơn lẻ hoặc văn bản tĩnh, bài toán này đòi hỏi mô hình phải có khả năng xử lý và tổng hợp tri thức từ các cấu trúc đồ thị đa quan hệ, đồng thời nắm bắt được tính sở hữu đặc thù để phân tích đúng bản chất của các cụm Sybil.\vspace{0.3cm}

Hệ thống phân lớp được xây dựng dựa trên ma trận đặc trưng nút có tổng cộng 396 chiều, kết hợp cùng đồ thị cấu trúc, cụ thể như sau:
\vspace{-0.1cm}
\begin{itemize}[label={--}, itemsep=1pt]
    \item \textbf{Đặc trưng ngữ nghĩa:} Khai thác thông tin từ hồ sơ cá nhân của người dùng bao gồm tên hiển thị, tên định danh, và tiểu sử. Dữ liệu văn bản này được nhúng qua mô hình ngôn ngữ Sentence-BERT (\texttt{all-MiniLM-L6-v2}) để tạo thành một vector 384 chiều, giúp mô hình nắm bắt được sự tương đồng về ngữ pháp, ngữ nghĩa và phát hiện các dấu vết rập khuôn, sao chép nội dung của các tài khoản ảo.
    
    \item \textbf{Đặc trưng thống kê và hành vi:} Bổ sung 12 chiều dữ liệu định lượng và trạng thái của tài khoản, bao gồm:
    \begin{itemize}[label={+}, itemsep=1pt]
        \item \textit{Chỉ số cường độ hoạt động:} 10 đặc trưng thống kê tổng quan (như tổng lượt theo dõi, đang theo dõi, bài đăng, bình luận, v.v.). Toàn bộ các chỉ số định lượng này đều được chuẩn hóa về khoảng $[0, 1]$ bằng \texttt{MinMaxScaler} để tránh mất cân bằng thang đo trong quá trình huấn luyện mạng nơ-ron.
        \item \textit{Trạng thái đại diện:} 1 đặc trưng nhị phân xác định tài khoản có sử dụng hình ảnh đại diện tùy chỉnh hay chỉ dùng ảnh mặc định.
        \item \textit{Độ tuổi tài khoản:} 1 đặc trưng tính toán số ngày hoạt động kể từ thời điểm khởi tạo.
    \end{itemize}
    
    \item \textbf{Đặc trưng cấu trúc mạng lưới:} Đây là luồng thông tin đóng vai trò then chốt, biểu diễn mối liên kết giữa các tài khoản dưới dạng đồ thị đa quan hệ có trọng số. Cấu trúc này không chỉ bao gồm các hành vi tương tác trên mạng xã hội (theo dõi và tương tác), khai thác tính sở hữu trong mạng xã hội Web3, và tầng tương đồng nhằm liên kết các tài khoản có sự rập khuôn về tiểu sử, tên định danh hoặc thời gian khởi tạo sát nhau. Sự kết hợp đa chiều này giúp mô hình bóc trần toàn diện cấu trúc và các thuật toán sinh tài khoản tự động của các mạng lưới Sybil.
\end{itemize}

Mục tiêu cốt lõi của bài toán là vượt qua giới hạn của việc thiếu nhãn dữ liệu thông qua việc học không giám sát (GAE kết hợp K-Means) và sau đó xây dựng một kiến trúc hợp nhất dựa trên Mạng nơ-ron đồ thị (GNN) kết hợp Máy học truyền thống (Decoupled ML). Kiến trúc này đóng vai trò nhúng các đặc trưng và cấu trúc mạng lưới vào một không gian chung, nhằm tối ưu hóa độ chính xác phân lớp và cung cấp khả năng giải thích.

%=======================================================================
\subsubsection{1.2. Ứng dụng Web (SybilSignal Dashboard)}
\indent\indent Để minh chứng cho tính thực tiễn của mô hình nghiên cứu và hỗ trợ người dùng cuối trong việc nhận diện rủi ro, đề tài triển khai xây dựng ứng dụng web \textbf{SybilSignal}. Hệ thống được thiết kế không chỉ để xuất ra kết quả dự đoán khô khan mà còn hướng tới việc trực quan hóa, giúp người dùng hiểu rõ nguyên nhân tại sao một tài khoản bị đánh giá là độc hại.\vspace{0.3cm}

Các chức năng chính của hệ thống bao gồm:\vspace{-0.1cm}
\begin{itemize}[label={--}, itemsep=1pt]
    \item \textbf{Chức năng 1 -- Tra cứu và dự đoán tức thời:} Người dùng nhập định danh (\texttt{profile\_id}) của một tài khoản cụ thể. Hệ thống sẽ kết nối với cơ sở dữ liệu để lấy dữ liệu và trích xuất đặc trưng, đi qua luồng suy luận của tổ hợp mô hình tối ưu (GNN + Machine Learning) và trả về mức độ rủi ro. Đi kèm với kết quả dự đoán của tài khoản mục tiêu, giao diện sẽ trực quan hóa đồ thị cục bộ thể hiện các kết nối của tài khoản đó với mạng lưới xung quanh từ đồ thị tham chiếu.
    \item \textbf{Chức năng 2 -- Thực thi Pipeline luồng thời gian:} Hệ thống cho phép người quản trị chọn một khoảng thời gian. Kích hoạt toàn bộ luồng tự động: từ việc thu thập dữ liệu, chạy Graph Autoencoder (GAE) để học cấu trúc và trích xuất đặc trưng đồ thị, dùng K-Means phân cụm, đến việc áp dụng hệ thống luật để tính điểm và dán nhãn. Sau đó, hệ thống sẽ hiển thị kết quả trực quan với đồ thị, hiển thị thông tin cụm và thông tin vi phạm của từng nút trong đồ thị.
\end{itemize}